\paragraph{Post-collection quality exclusion.}
The registered protocol specified no annotator exclusions. After inspecting the
first four completed packets, we identified one reviewer who selected the same
directional response on 17 of 18 items despite the packet's balanced registered
directions (six pro-YES, six anti-YES, and six null). We classified this pattern
as failure to perform the conditional-direction task and adopted a uniform
post-hoc rule excluding any completed reviewer with at least 17 identical
directional responses. The rule excluded one of seven recruited reviewers. We
retained all raw responses, analyzed the six remaining reviewers with a
five-of-six item gate, and report the original seven-reviewer analysis as a
sensitivity check. No item was removed or relabeled.

After the rule was adopted, and after R2 had been told their score and response
pattern, a follow-up debrief provided context for the failure. R2 reported
interpreting the task as asking whether reports supported YES, difficulty with
English and market-specific terminology, consulting a dictionary and an AI
assistant despite the no-outside-assistance instruction, and reduced attention.
No re-rating was accepted. Because this debrief followed performance disclosure
and researcher-provided examples, we treat it as contextual information rather
than independent evidence for exclusion. The exclusion rests only on the uniform
response-pattern rule.
